\documentclass[12pt]{article}

% Ustawienia strony
\usepackage{lipsum}
\usepackage[margin=1.5cm]{geometry}
\usepackage{graphicx}
\usepackage{enumitem}
\usepackage{siunitx}
\usepackage{amsmath}
\usepackage{cancel}
\usepackage{makecell}
\usepackage{mathtools}
\usepackage{floatrow}
\usepackage{caption}
\usepackage[T1]{fontenc}
\usepackage{fancyvrb}
\usepackage{subcaption}
\usepackage{blindtext}
\usepackage{hyperref}
\usepackage{amssymb}
\usepackage{minted}
\usepackage{xcolor} % to access the named color LightGray
\definecolor{LightGray}{gray}{0.9}

\setminted{
	frame=lines,
	framesep=1mm,
	baselinestretch=1,
	bgcolor=LightGray,
	fontsize=\footnotesize,
	linenos
}

\hypersetup{
	colorlinks=true,
	linkcolor=blue,
	filecolor=magenta,      
	urlcolor=cyan,
	pdfpagemode=FullScreen,
}

\urlstyle{same}

\fvset{tabsize=4}
\setlist{leftmargin=15pt,labelindent=15pt}
\setlist[enumerate]{noitemsep,wide=0pt, leftmargin=*}
\setlength{\parindent}{0pt}

% Ustawienia stopki
\usepackage{lastpage}
\usepackage{fancyhdr}
\pagestyle{fancy}
\fancyhf{} % sets both header and footer to nothing
\renewcommand{\headrulewidth}{0pt}
\cfoot{\thepage\ z \pageref*{LastPage}}

\begin{document}
\textbf{Robert Mesek}\text{,} % Autor
\text{05.06.2025}\text{ } % Data

\begin{center}
  \section*{Optymalizacja Kodu na Różne Architektury \\
    Zadanie 2 - Eliminacja Gaussa}
  \mbox{}
\end{center}

\section*{Cel zadania}
Celem zadania jest implementacja i optymalizacja algorytmu eliminacji Gaussa,
z uwzględnieniem mikroarchitektury posiadanego procesora oraz praktycznego
wykorzystania technik niskopoziomowej optymalizacji kodu. \\ \\
Zakres prac:
\begin{itemize}
  \item Zaimplementuj algorytm eliminacji Gaussa (z lub bez pivotingu).
  \item Zoptymalizuj implementację pod kątem architektury CPU (cache,
        SIMD, itp.).
  \item Zweryfikuj poprawność: wynik optymalizowanej wersji musi odpowiadać
        wersji referencyjnej.
  \item Oszacuj liczby operacji (FLOP) i przelicz GFLOPS.
  \item Przeprowadź pomiary wydajności (opcjonalnie z użyciem PAPI lub
        równoważnego narzędzia).
\end{itemize}
Wymagania techniczne:
\begin{itemize}
  \item Wektoryzacja obliczeń z użyciem instrukcji SIMD (np. AVX, SSE).
  \item Dbałość o lokalność danych (wykorzystanie pamięci cache).
  \item Analiza architektury CPU i dostosowanie kodu do jednostek wykonawczych.
  \item Pomiary czasu wykonania oraz obliczonych GFLOPS-ów.
  \item Wykonanie testów poprawności dla różnych rozmiarów macierzy.
\end{itemize}


\newpage

\section*{Dane techniczne}
\subsection*{System Operacyjny}
\begin{minted}{text}
$ hostnamectl
Virtualization: wsl
Operating System: Ubuntu 24.04.2 LTS
Kernel: Linux 6.6.87.1-microsoft-standard-WSL2
Architecture: x86-64
\end{minted}

\subsection*{Kompilator}
\begin{minted}{text}
$ gcc --version
gcc (Ubuntu 13.3.0-6ubuntu2~24.04) 13.3.0
\end{minted}

\subsection*{Procesor}
\begin{minted}{text}
$ lscpu
Architecture:             x86_64
  CPU op-mode(s):         32-bit, 64-bit
  Address sizes:          48 bits physical, 48 bits virtual
  Byte Order:             Little Endian
CPU(s):                   24
  On-line CPU(s) list:    0-23
Vendor ID:                AuthenticAMD
  Model name:             AMD Ryzen 9 7900X 12-Core Processor
    CPU family:           25
    Model:                97
    Thread(s) per core:   2
    Core(s) per socket:   12
    Socket(s):            1
    Stepping:             2
    BogoMIPS:             9399.78
    Flags:                fpu vme de pse tsc msr pae mce cx8 apic sep mtrr pge mca cmov pat pse36 
                          clflush mmx fxsr sse sse2 ht syscall nx mmxext fxsr_opt pdpe1gb rdtscp 
                          lm constant_tsc rep_good nopl tsc_reliable nonsto p_tsc cpuid extd_apicid 
                          pni pclmulqdq ssse3 fma cx16 sse4_1 sse4_2 movbe popcnt aes xsave avx f16c 
                          rdrand hypervisor lahf_lm cmp_legacy svm cr8_legacy abm sse4a misalignsse 
                          3dnowprefetch osvw topoext perfctr_core ssbd ibrs ibpb stibp vmmcall 
                          fsgsbase bmi1 avx2 smep bmi2 erms invpcid avx512f avx512dq rdseed adx smap 
                          avx512ifma clflushopt clwb avx512cd sha_ni avx512bw avx512vl xsaveo pt 
                          xsavec xgetbv1 xsaves avx512_bf16 clzero xsaveerptr arat npt nrip_save 
                          tsc_scale vmcb_clean flushbyasid decodeassists pausefilter pfthreshold 
                          v_vmsave_vmload avx512vbmi umip avx512_vbmi2 gf ni vaes vpclmulqdq 
                          avx512_vnni avx512_bitalg avx512_vpopcntdq rdpid fsrm
Caches (sum of all):      
  L1d:                    384 KiB (12 instances)
  L1i:                    384 KiB (12 instances)
  L2:                     12 MiB (12 instances)
  L3:                     32 MiB (1 instance)
\end{minted}

\newpage














\section*{Weryfikacja poprawności}
Układy testowych danych zostały wygenerowane przy pomocy funkcji
\begin{minted}{c}
void generate_solvable_system(double **A, double *b, double *true_x,
                              const int SIZE) {
  for (int i = 0; i < SIZE; i++) {
    true_x[i] = ((double)rand() / RAND_MAX) * 20.0 - 10.0;
  }

  for (int i = 0; i < SIZE; i++) {
    for (int j = 0; j < SIZE; j++) {
      A[i][j] = ((double)rand() / RAND_MAX) * 20.0 - 10.0;
    }
  }

  for (int i = 0; i < SIZE; i++) {
    b[i] = 0.0;
    for (int j = 0; j < SIZE; j++) {
      b[i] += A[i][j] * true_x[j];
    }
  }
}
\end{minted}

Sprawdzenie rozwiązania wykonano przy pomocy funkcji
\begin{minted}{c}
int verify_solution(double *x, double *true_x, const int SIZE) {
  double tolerance = 1e-4;

  for (int i = 0; i < SIZE; i++) {
    if (fabs(x[i] - true_x[i]) > tolerance) return 1;  // Verification failed
  }
  return 0;  // Verification successful
}
\end{minted}

Zaimplementowana wersja algorytmu \textbf{nie zawierała pivotingu}, a testy były przeprowadzane na \\
\textbf{macierzach kwadratowych} \textbf{bez wielowątkowości}.

\section*{Obliczenie GFLOPS-ów}
Do obliczenia GFLOPS-ów wykorzystano narzędzie \verb|perf|.
Przykładowy pomiar
\begin{minted}{text}
$ gcc -mfma -O2 ge3.c -o ge3
$ perf stat -e fp_ret_sse_avx_ops.all,task-clock ./ge3 3333
Time (seconds): 1.851049e+00

 Performance counter stats for './ge3 3333':

       24767274143      fp_ret_sse_avx_ops.all:u         #   13.920 G/sec
           1779.26 msec task-clock:u                     #    0.908 CPUs utilized

       1.958770318 seconds time elapsed

       1.729572000 seconds user
       0.047236000 seconds sys
\end{minted}

\newpage








\section*{Wersja podstawowa}
Algorytm eliminacji Gaussa bez optimalizacji. 
\begin{minted}{c}
int gaussian_elimination(double **A, double *b, double *x, const int SIZE) {
  // Forward elimination
  for (int k = 0; k < SIZE - 1; k++) {
    // Eliminate column k in rows below
    for (int i = k + 1; i < SIZE; i++) {
      double factor = A[i][k] / A[k][k];
      for (int j = k; j < SIZE; j++) {
        A[i][j] -= factor * A[k][j];
      }
      b[i] -= factor * b[k];
    }
  }

  // Back substitution
  for (int i = SIZE - 1; i >= 0; i--) {
    x[i] = b[i];
    for (int j = i + 1; j < SIZE; j++) {
      x[i] -= A[i][j] * x[j];
    }
    x[i] /= A[i][i];
  }

  return 0;  // Success
}
\end{minted}

\begin{table}[H]
  \centering
  \begin{tabular}{|c|c|c|c|}
    \hline
    \textbf{Macierz (NxN)} & \textbf{Czas algorytmu (s)} & \textbf{Całkowity czas (s)} & \textbf{Całkowite GFLOPS} \\
    \hline
    \verb|333|  & \verb|2.709000e-03| & \verb|4.723702e-03| & \verb|5.956| \\
    \hline
    \verb|1111| & \verb|1.090970e-01| & \verb|1.210884e-01| & \verb|7.830| \\
    \hline
    \verb|2222| & \verb|1.215120e+00| & \verb|1.265453e+00| & \verb|5.853| \\
    \hline
    \verb|3333| & \verb|6.197356e+00| & \verb|6.328828e+00| & \verb|4.265| \\
    \hline
    \verb|6666| & \verb|4.859530e+01| & \verb|5.015407e+01| & \verb|4.009| \\
    \hline
  \end{tabular}
  \caption{Wyniki pomiarów (perf) dla wersji podstawowej}
\end{table}

\newpage







\section*{Wersja z wektoryzacją}
Algorytm eliminacji Gaussa z wektoryzacją przy użyciu instrukcji AVX2. 
\begin{minted}{c}
void kernel_eliminate_8(double *A_pivot, double *A_target, int jj,
                        double factor, const int SIZE) {
  register __m256d factor_vec;
  register __m256d pivot_vec_1, pivot_vec_2;
  register __m256d target_vec_1, target_vec_2;

  factor_vec = _mm256_set1_pd(factor);

  pivot_vec_1 = _mm256_loadu_pd(&A_pivot[jj]);
  pivot_vec_2 = _mm256_loadu_pd(&A_pivot[jj + 4]);

  target_vec_1 = _mm256_loadu_pd(&A_target[jj]);
  target_vec_2 = _mm256_loadu_pd(&A_target[jj + 4]);

  target_vec_1 = _mm256_fnmadd_pd(factor_vec, pivot_vec_1, target_vec_1);
  target_vec_2 = _mm256_fnmadd_pd(factor_vec, pivot_vec_2, target_vec_2);

  _mm256_storeu_pd(&A_target[jj], target_vec_1);
  _mm256_storeu_pd(&A_target[jj + 4], target_vec_2);
}

int gaussian_elimination(double **A, double *b, double *x, const int SIZE) {
  double factor;
  int i, j, k, jj;

  // Forward elimination
  for (k = 0; k < SIZE - 1; k++) {
    for (i = k + 1; i < SIZE; i++) {
      factor = A[i][k] / A[k][k];
      for (jj = k; jj <= SIZE - 8; jj += 8) {
        kernel_eliminate_8(A[k], A[i], jj, factor, SIZE);
      }
      for (j = jj; j < SIZE; j++) A[i][j] -= factor * A[k][j];
      b[i] -= factor * b[k];
    }
  }

  // Back substitution
  for (int i = SIZE - 1; i >= 0; i--) {
    x[i] = b[i];
    for (int j = i + 1; j < SIZE; j++) {
      x[i] -= A[i][j] * x[j];
    }
    x[i] /= A[i][i];
  }

  return 0;  // Success
}
\end{minted}

\begin{table}[H]
  \centering
  \begin{tabular}{|c|c|c|c|}
    \hline
    \textbf{Macierz (NxN)} & \textbf{Czas algorytmu (s)} & \textbf{Całkowity czas (s)} & \textbf{Całkowite GFLOPS} \\
    \hline
    \verb|333|  & \verb|1.629000e-03| & \verb|3.580134e-03| & \verb|8.041| \\
    \hline
    \verb|1111| & \verb|6.874200e-02| & \verb|8.100482e-02| & \verb|11.642| \\
    \hline
    \verb|2222| & \verb|8.545080e-01| & \verb|9.089234e-01| & \verb|8.186| \\
    \hline
    \verb|3333| & \verb|4.543042e+00| & \verb|4.670241e+00| & \verb|5.718| \\
    \hline
    \verb|6666| & \verb|3.929290e+01| & \verb|4.025762e+01| & \verb|4.989| \\
    \hline
  \end{tabular}
  \caption{Wyniki pomiarów (perf) dla wersji z wektoryzacją}
\end{table}

\newpage






\section*{Wersja z wektoryzacją i blokowaniem}
Algorytm eliminacji Gaussa z wektoryzacją przy użyciu instrukcji AVX2 oraz efektywniejszym wykorzystaniem pamięci cache poprzez blokowanie.
\begin{minted}{c}
#define BLOCK_SIZE 32

void kernel_eliminate_8(double *A_pivot, double *A_target, int jj,
                        double factor, const int SIZE) {
  register __m256d factor_vec;
  register __m256d pivot_vec_1, pivot_vec_2;
  register __m256d target_vec_1, target_vec_2;

  factor_vec = _mm256_set1_pd(factor);

  pivot_vec_1 = _mm256_loadu_pd(&A_pivot[jj]);
  pivot_vec_2 = _mm256_loadu_pd(&A_pivot[jj + 4]);

  target_vec_1 = _mm256_loadu_pd(&A_target[jj]);
  target_vec_2 = _mm256_loadu_pd(&A_target[jj + 4]);

  target_vec_1 = _mm256_fnmadd_pd(factor_vec, pivot_vec_1, target_vec_1);
  target_vec_2 = _mm256_fnmadd_pd(factor_vec, pivot_vec_2, target_vec_2);

  _mm256_storeu_pd(&A_target[jj], target_vec_1);
  _mm256_storeu_pd(&A_target[jj + 4], target_vec_2);
}

int gaussian_elimination(double **A, double *b, double *x, const int SIZE) {
  double factor;
  int i, j, k, jj, kk, kk_end;

  // Forward elimination
  for (kk = 0; kk < SIZE - 1; kk += BLOCK_SIZE) {
    kk_end = (kk + BLOCK_SIZE < SIZE) ? kk + BLOCK_SIZE : SIZE;

    // Process current block
    for (k = kk; k < kk_end - 1; k++) {
      for (i = k + 1; i < kk_end; i++) {
        factor = A[i][k] / A[k][k];
        for (jj = k; jj <= SIZE - 8; jj += 8) {
          kernel_eliminate_8(A[k], A[i], jj, factor, SIZE);
        }
        for (j = jj; j < SIZE; j++) A[i][j] -= factor * A[k][j];
        b[i] -= factor * b[k];
      }
    }

    // Update remaining rows using the processed block
    for (i = kk_end; i < SIZE; i++) {
      for (k = kk; k < kk_end; k++) {
        factor = A[i][k] / A[k][k];
        for (jj = k; jj <= SIZE - 8; jj += 8) {
          kernel_eliminate_8(A[k], A[i], jj, factor, SIZE);
        }
        for (j = jj; j < SIZE; j++) A[i][j] -= factor * A[k][j];
        b[i] -= factor * b[k];
      }
    }
  }

  // Back substitution
  for (int i = SIZE - 1; i >= 0; i--) {
    x[i] = b[i];
    for (int j = i + 1; j < SIZE; j++) {
      x[i] -= A[i][j] * x[j];
    }
    x[i] /= A[i][i];
  }

  return 0;  // Success
}
\end{minted}

\begin{table}[H]
  \centering
  \begin{tabular}{|c|c|c|c|}
    \hline
    \textbf{Macierz (NxN)} & \textbf{Czas algorytmu (s)} & \textbf{Całkowity czas (s)} & \textbf{Całkowite GFLOPS} \\
    \hline
    \verb|333|  & \verb|1.447000e-03| & \verb|3.285216e-03| & \verb|8.875| \\
    \hline
    \verb|1111| & \verb|5.827000e-02| & \verb|7.104668e-02| & \verb|13.227| \\
    \hline
    \verb|2222| & \verb|4.786020e-01| & \verb|5.327180e-01| & \verb|13.986| \\
    \hline
    \verb|3333| & \verb|1.851049e+00| & \verb|1.958770e+00| & \verb|13.920| \\
    \hline
    \verb|6666| & \verb|1.482420e+01| & \verb|1.588368e+01| & \verb|12.712| \\
    \hline
  \end{tabular}
  \caption{Wyniki pomiarów (perf) dla wersji z wektoryzacją i blokowaniem}
\end{table}

\newpage








\section*{Wnioski}
Czas oraz liczba operacji w generowaniu macierzy, weryfikacji rozwiązania 
oraz w samym algorytmie eliminacji podczas podstawienia wstecznego były pomijalnie
małe w porównaniu do czasu eliminacji, więc do obliczenia GFLOPS-ów wystarczyło użycie
całkowitego czasu wykonania programu. \\

Prefetchowanie danych mierzalnie nie wpływało na czas wykonania lub nawet pogarszało 
wyniki dla wersji z wektoryzacją i blokowaniem, więc zostało pominięte. \\

Całkowita liczba operacji zmiennoprzecinkowych przybliżona jest do teoretycznej złożoności
$(2/3)N^3$, przy macierzy $N \times N$. \\

Zastosowanie optymalizacji znacznie poprawiło wyniki algorytmu. Poprawa zależy od rozmiaru macierzy,
gdzie dla większych macierzy różnice są bardziej zauważalne. Ręczna implementacja wektoryzacji 
przyniosła poprawę o około 25\% przy rozmiarze macierzy 3333, co może wynikać z faktu, że 
kompilator tylko w pewnym stopniu wykorzystał automatyczną wektoryzację.
Wersja z blokowaniem natomiast przyniosła dużą poprawę o około 50\% przy rozmiarze macierzy 3333, 
więc widzimy, że w tym przypadku efektywne ładowanie danych do pamięci cache ma kluczowe znaczenie. 

\end{document}
